
\clearpage
\phantomsection
\addcontentsline{toc}{chapter}{CHƯƠNG 4: THIẾT KẾ VẬT LÝ}
{\LARGE\bfseries CHƯƠNG 4: THIẾT KẾ VẬT LÝ}\par\vspace{1em}

\phantomsection
\addcontentsline{toc}{section}{4.1. Lựa chọn hệ quản trị CSDL}
{\Large\bfseries 4.1. Lựa chọn hệ quản trị CSDL}\par\vspace{0.7em}

Đồ án triển khai vật lý trên Microsoft SQL Server 2022 (chạy trong container Docker). Việc chọn SQL Server thay vì PostgreSQL/MySQL bảo đảm nhất quán với môi trường thực hành đã cài đặt trong học phần và cho phép sinh viên kiểm chứng trực tiếp bằng sqlcmd hoặc Azure Data Studio.

\phantomsection
\addcontentsline{toc}{section}{4.2. Ánh xạ kiểu dữ liệu logic sang vật lý}
{\Large\bfseries 4.2. Ánh xạ kiểu dữ liệu logic sang vật lý}\par\vspace{0.7em}

\begin{longtable}{|p{0.3\textwidth}|p{0.3\textwidth}|p{0.3\textwidth}|}
\hline
\textbf{Khái niệm logic} & \textbf{Kiểu vật lý SQL Server} & \textbf{Ghi chú} \\ \hline
\endfirsthead
\hline
\textbf{Khái niệm logic} & \textbf{Kiểu vật lý SQL Server} & \textbf{Ghi chú} \\ \hline
\endhead
Khóa thay thế (UUID) & UNIQUEIDENTIFIER DEFAULT NEWID() & Sinh tự động, thuận tiện đồng bộ dữ liệu giữa Web/Mobile/AI Service \\ \hline
Văn bản tiếng Việt có dấu & NVARCHAR(n) / NVARCHAR(MAX) & Bắt buộc dùng NVARCHAR để lưu Unicode đúng, tránh mất dấu \\ \hline
Mã, URL, email (ASCII) & VARCHAR(n) & Tiết kiệm không gian lưu trữ so với NVARCHAR \\ \hline
Thời điểm & DATETIME2 & Độ chính xác cao hơn DATETIME cũ; lưu theo UTC, quy ước ở tầng ứng dụng \\ \hline
Cờ đúng/sai & BIT & Tương đương BOOLEAN \\ \hline
Dữ liệu bán cấu trúc (JSON) & NVARCHAR(MAX) + CHECK (ISJSON(col)=1) & SQL Server không có kiểu JSONB như PostgreSQL; dùng hàm ISJSON để ràng buộc hợp lệ \\ \hline
Vector embedding AI & NVARCHAR(MAX) (mảng số dạng JSON) & SQL Server 2022 chưa có kiểu VECTOR gốc (khác PostgreSQL/pgvector); có thể tích hợp thêm Azure AI Search hoặc vector store ngoài nếu cần tìm kiếm gần đúng (ANN) hiệu năng cao \\ \hline
Số tự tăng (log, cấu hình) & INT/BIGINT IDENTITY(1,1) & Dùng cho bảng log tần suất ghi cao, không cần tính ngẫu nhiên của UUID \\ \hline
\end{longtable}


\phantomsection
\addcontentsline{toc}{section}{4.3. Ràng buộc toàn vẹn và chính sách xóa/cập nhật}
{\Large\bfseries 4.3. Ràng buộc toàn vẹn và chính sách xóa/cập nhật}\par\vspace{0.7em}

Theo mục 'Chính sách cập nhật và xóa' của Chương 2, mỗi khóa ngoại được gán chính sách phù hợp ngữ nghĩa:

\begin{itemize}
\item CASCADE: xóa GIA\_DINH sẽ xóa lan truyền CHI\_NHANH, BAI\_DANG, SU\_KIEN, tài liệu di sản... vì các đối tượng này không có ý nghĩa khi dòng họ không còn tồn tại.
\item NO ACTION (RESTRICT thủ công ở tầng ứng dụng): FamilyMembers.BranchId và EventParticipants.MemberId dùng NO ACTION để tránh lỗi 'multiple cascade paths' của SQL Server khi một bảng nhận nhiều đường xóa lan truyền từ cùng gốc GIA\_DINH; việc dọn dữ liệu liên quan được xử lý bằng thủ tục ứng dụng hoặc trigger AFTER DELETE.
\item Thực thể yếu (UserProfiles, MemberVerifications, EventReminders) đều dùng CASCADE, đúng khuyến nghị Chương 2: 'với thực thể yếu, CASCADE thường phù hợp vì phụ thuộc vòng đời'.
\end{itemize}

\phantomsection
\addcontentsline{toc}{section}{4.4. Chỉ mục (Index)}
{\Large\bfseries 4.4. Chỉ mục (Index)}\par\vspace{0.7em}

Chỉ mục được thiết kế theo khối lượng truy vấn dự kiến: liệt kê thành viên/chi nhánh theo dòng họ, bài đăng/sự kiện theo dòng họ, bình luận theo bài đăng, và tra cứu embedding AI theo (EntityType, EntityId). Danh sách đầy đủ nằm ở cuối script Phụ lục A.

\phantomsection
\addcontentsline{toc}{section}{4.5. Trích đoạn script T-SQL}
{\Large\bfseries 4.5. Trích đoạn script T-SQL}\par\vspace{0.7em}

Toàn bộ script tạo CSDL FamilyConnectDB được trình bày đầy đủ ở Phụ lục A và đính kèm dưới dạng tệp schema\_sqlserver.sql. Dưới đây là trích đoạn minh họa nhóm Sự kiện — đúng kết quả suy ra từ phân tích chuẩn hóa ở Chương 3:

\begin{lstlisting}[language=SQL]
CREATE TABLE Events (
    EventId         UNIQUEIDENTIFIER NOT NULL DEFAULT NEWID() PRIMARY KEY,
    FamilyId        UNIQUEIDENTIFIER NOT NULL REFERENCES Families(FamilyId) ON DELETE CASCADE,
    CreatedBy       UNIQUEIDENTIFIER NULL REFERENCES Users(UserId),
    Title           NVARCHAR(200) NOT NULL,
    StartTime       DATETIME2 NOT NULL,
    EndTime         DATETIME2 NULL,
    ...
);
 
CREATE TABLE EventParticipants (
    EventId         UNIQUEIDENTIFIER NOT NULL REFERENCES Events(EventId) ON DELETE CASCADE,
    MemberId        UNIQUEIDENTIFIER NOT NULL REFERENCES FamilyMembers(MemberId) ON DELETE NO ACTION,
    RsvpStatus      VARCHAR(20) NOT NULL DEFAULT 'invited'
                        CHECK (RsvpStatus IN ('invited','accepted','declined','maybe')),
    PRIMARY KEY (EventId, MemberId)
);
 
-- Thực thể yếu: khóa chủ + khóa cục bộ (RemindAt), đúng Thuật toán 2.3
CREATE TABLE EventReminders (
    EventId         UNIQUEIDENTIFIER NOT NULL REFERENCES Events(EventId) ON DELETE CASCADE,
    RemindAt        DATETIME2 NOT NULL,
    Channel         VARCHAR(20) NOT NULL DEFAULT 'push',
    PRIMARY KEY (EventId, RemindAt)
);
\end{lstlisting}

\phantomsection
\addcontentsline{toc}{section}{4.6. Kiểm thử bằng dữ liệu mẫu}
{\Large\bfseries 4.6. Kiểm thử bằng dữ liệu mẫu}\par\vspace{0.7em}

Theo khuyến nghị 'Kiểm chứng bằng dữ liệu mẫu' của Chương 2, các tình huống biên sau được dùng để kiểm thử lược đồ vật lý sau khi triển khai:

\begin{itemize}
\item Thêm một Family mới chưa có FamilyMember nào — hợp lệ vì Families không phụ thuộc FamilyMembers.
\item Thêm một FamilyMember chưa có UserId liên kết (thành viên đã mất, không có tài khoản) — hợp lệ vì UserId cho phép NULL.
\item Thêm hai FamilyMember cùng ParentMemberId khác nhau cho một ChildMemberId (hai cha/mẹ) — hợp lệ với 2 dòng trong ParentChildRelationships.
\item Xóa một Event còn EventParticipants — CASCADE tự xóa toàn bộ dòng liên quan trong EventParticipants và EventReminders.
\item Thử thêm EventReminders trùng (EventId, RemindAt) — bị chặn bởi khóa chính ghép, đúng ngữ nghĩa khóa cục bộ của thực thể yếu.
\end{itemize}